\section{Introducción}

El presente informe corresponde al proyecto final del curso Construcción y Sustentabilidad en Hormigón Armado, en el cual se aplican los contenidos revisados a lo largo del semestre a un caso real de construcción. El proyecto consiste en la nueva Planta de Tratamiento de Aguas Servidas para Reúso de Antofagasta, emplazada en el sector industrial de Salar del Carmen--La Negra, alejada del centro urbano de la ciudad. Entre las principales estructuras de la planta se consideran los reactores biológicos, los clarificadores secundarios de forma circular y un diámetro interior de 26 metros, y el campamento de los trabajadores.

El emplazamiento condiciona de manera transversal todas las decisiones del proyecto. El sector se ubica en pleno Desierto de Atacama, en un clima desértico árido extremo caracterizado por muy baja humedad relativa, alta radiación solar y vientos sostenidos, además de una ubicación remota que encarece la logística de materiales y de mano de obra. A esto se suma que, por tratarse de estructuras de contención de líquidos en contacto con aguas residuales, los clarificadores y reactores imponen exigencias estrictas de impermeabilidad y durabilidad. Estos factores, junto con una duración de obra estimada en al menos dos años, obligan a considerar el hormigonado en condiciones extremas tanto de verano como de invierno.

El objetivo de este trabajo es desarrollar, para una de las estructuras de tratamiento y para el campamento, una propuesta técnica que abarque su proceso constructivo, la especificación de los hormigones según sus condiciones de exposición, la identificación y reparación de los problemas de calidad más probables, y la incorporación de un método moderno de construcción. Para ello, el informe se estructura en cuatro secciones, alineadas con los requerimientos del enunciado. La primera describe el proceso constructivo del clarificador secundario y del campamento, incluyendo moldajes, colocación, compactación, curado y juntas, así como las consideraciones de hormigonado en condiciones extremas. La segunda aborda la especificación del hormigón conforme a la NCh170:2016 y plantea el uso de un hormigón especial con su respectivo análisis de costos. La tercera identifica los principales problemas de calidad y su metodología de reparación. La cuarta propone la aplicación de un método moderno de construcción a partir de lo observado en la visita al Parque CTEC.

Todos los supuestos adoptados en el desarrollo del proyecto se indican de forma explícita y debidamente justificada en cada sección, dado que el enunciado no especifica la totalidad de los parámetros de diseño de las estructuras.